%!TEX root = ../main.tex

\chapter{Implementierung}
Auf Basis der in Kapitel \ref{chap:50_AuswahlImplementierungsansatz} beschriebenen Auswahl des Implementierungsansatzes wird dieser in 
diesem Kapitel umgesetzt. Das Motion Planning besteht aus zwei Schritten: der Berechnung der Delaunay Triangulation und der Filterung
der Trackdreiecke. Und im zweiten Schritt aus der Berechnung der Centerline. Hier wird ein Tree-Search-Algorithmus eingesetzt, um aus den 
gegebenen Trackdreiecken, die vorteilhafteste Route zu berechnen. Abbildung \ref{fig:ÜbersichtImplementierung} zeigt eine Übersicht des 
implementierten Systems. Auf die Schnittstellen, die einzelnen Module und deren Funktionen wird im Folgenden näher eingegangen. 

\subsection{Schnittstellen}
Als Steuerungseingänge werden die in Kapitel \ref{chap:20_Grundlagen} beschriebenen Steuersignale abboniert...

Als primäre Signaleingänge werden die erkannten Pylonen und die aktuelle Fahrzeugposition genutzt. Beide dieser 

\subsection{Delaunay Triangulation}

\subsection{Filterung von Trackdreiecken}

\subsection{Centerline}

